\documentclass[UTF8,aspectratio=169]{ctexbeamer}

\usetheme{Madrid}
\usecolortheme{default}
\setbeamertemplate{navigation symbols}{}

\usepackage{booktabs}
\usepackage{tabularx}
\usepackage{array}
\usepackage{ragged2e}
\usepackage{hyperref}
\usepackage{amsmath}
\usepackage{tikz}

\hypersetup{hidelinks}
\urlstyle{tt}

\definecolor{mmBlue}{HTML}{123B5D}
\definecolor{mmTeal}{HTML}{1E6F78}
\definecolor{mmGold}{HTML}{C98B2E}
\definecolor{mmGray}{HTML}{4A5568}
\definecolor{mmLite}{HTML}{EEF3F7}

\setbeamercolor{title}{fg=white,bg=mmBlue}
\setbeamercolor{frametitle}{fg=white,bg=mmBlue}
\setbeamercolor{structure}{fg=mmTeal}
\setbeamercolor{block title}{fg=white,bg=mmTeal}
\setbeamercolor{block body}{fg=black,bg=mmLite}
\setbeamercolor{footline}{fg=white,bg=mmBlue}

\setbeamertemplate{footline}{
  \leavevmode%
  \hbox{%
    \begin{beamercolorbox}[wd=.82\paperwidth,ht=2.5ex,dp=1ex,leftskip=1em]{footline}%
      Linux mm 新手学习路线与版本演进
    \end{beamercolorbox}%
    \begin{beamercolorbox}[wd=.18\paperwidth,ht=2.5ex,dp=1ex,center]{footline}%
      \insertframenumber/\inserttotalframenumber
    \end{beamercolorbox}%
  }
}

\newcolumntype{Y}{>{\RaggedRight\arraybackslash}X}
\newcommand{\kpath}[1]{\path{#1}}
\newcommand{\kcode}[1]{\path{#1}}

\title{Linux \texttt{mm} 新手学习路线图}
\subtitle{从历史版本到当前主线（基于当前源码树 v7.0.0-rc7）}
\author{Codex 生成}
\date{2026-04-11}

\begin{document}

\begin{frame}
  \titlepage
\end{frame}

\begin{frame}[t]{这份 PPT 的定位}
\small
\begin{block}{目标}
把一个对 Linux 内核 \texttt{mm} 子系统几乎零基础的人，带到能够：
\begin{itemize}
  \item 看懂当前主线 \texttt{mm/}、\texttt{Documentation/mm/} 的主要结构。
  \item 读旧资料时知道哪些概念还有效，哪些术语已经换代。
  \item 能沿着“分配、缺页、页缓存、回收、控制”五条主线跟代码。
  \item 具备继续深入某个方向并开始提交小补丁的能力。
\end{itemize}
\end{block}

\begin{block}{范围}
\begin{itemize}
  \item 主线聚焦 MMU 系统；\texttt{nommu} 只做背景说明，不作为新手主线。
  \item 以当前源码树 \texttt{7.0.0-rc7} 为“现在”的参考点。
  \item 版本演进强调“理解结构变化的主线”，不是穷举每一个历史提交。
\end{itemize}
\end{block}
\end{frame}

\begin{frame}[t]{先把 \texttt{mm} 看成四层地图}
\small
\begin{columns}[T]
  \begin{column}{0.48\textwidth}
    \begin{block}{第 1 层：物理内存}
    \begin{itemize}
      \item node / zone / page frame / folio
      \item buddy allocator
      \item per-cpu pageset
      \item compaction / migration / hotplug
    \end{itemize}
    \end{block}

    \begin{block}{第 2 层：虚拟地址空间}
    \begin{itemize}
      \item \texttt{mm\_struct}
      \item \texttt{vm\_area\_struct}
      \item page tables
      \item page fault / COW / GUP / rmap
    \end{itemize}
    \end{block}
  \end{column}
  \begin{column}{0.48\textwidth}
    \begin{block}{第 3 层：文件与缓存}
    \begin{itemize}
      \item page cache / readahead
      \item writeback / dirty throttling
      \item shmem / tmpfs
      \item folio 化后的 file mm 路径
    \end{itemize}
    \end{block}

    \begin{block}{第 4 层：压力与控制}
    \begin{itemize}
      \item reclaim / swap / OOM
      \item memcg / NUMA policy
      \item THP / hugetlb / multigen LRU
      \item DAMON / HMM / memory tiers
    \end{itemize}
    \end{block}
  \end{column}
\end{columns}

\vspace{0.4em}
\small
学习顺序建议永远是：先地图，再主路径，再高级特性，再性能与调优。
\end{frame}

\begin{frame}[t]{当前主线源码树里应该先记住的入口}
\footnotesize
\begin{tabularx}{\textwidth}{p{0.20\textwidth}YY}
\toprule
主题 & 关键文档 & 关键代码 \\
\midrule
总体入口 &
\kpath{Documentation/mm/index.rst} \newline
\kpath{Documentation/admin-guide/mm/index.rst} &
\kpath{mm/} \newline
\kpath{include/linux/mm_types.h} \newline
\kpath{include/linux/mmzone.h} \\
\addlinespace
物理内存与分配 &
\kpath{Documentation/mm/physical_memory.rst} \newline
\kpath{Documentation/mm/page_allocation.rst} &
\kpath{mm/page_alloc.c} \newline
\kpath{mm/compaction.c} \newline
\kpath{mm/mm_init.c} \newline
\kpath{mm/memblock.c} \\
\addlinespace
地址空间与缺页 &
\kpath{Documentation/mm/process_addrs.rst} \newline
\kpath{Documentation/mm/page_tables.rst} &
\kpath{mm/mmap.c} \newline
\kpath{mm/vma.c} \newline
\kpath{mm/memory.c} \newline
\kpath{mm/rmap.c} \\
\addlinespace
页缓存与文件路径 &
\kpath{Documentation/mm/page_cache.rst} \newline
\kpath{Documentation/mm/vmalloc.rst} &
\kpath{mm/filemap.c} \newline
\kpath{mm/readahead.c} \newline
\kpath{mm/truncate.c} \newline
\kpath{mm/shmem.c} \\
\addlinespace
回收、交换、OOM &
\kpath{Documentation/mm/page_reclaim.rst} \newline
\kpath{Documentation/mm/swap.rst} \newline
\kpath{Documentation/mm/oom.rst} &
\kpath{mm/vmscan.c} \newline
\kpath{mm/workingset.c} \newline
\kpath{mm/swap.c} \newline
\kpath{mm/swapfile.c} \newline
\kpath{mm/oom_kill.c} \\
\addlinespace
现代主题 &
\kpath{Documentation/mm/transhuge.rst} \newline
\kpath{Documentation/mm/multigen_lru.rst} \newline
\kpath{Documentation/mm/damon/design.rst} \newline
\kpath{Documentation/mm/hmm.rst} &
\kpath{mm/huge_memory.c} \newline
\kpath{mm/khugepaged.c} \newline
\kpath{mm/memcontrol.c} \newline
\kpath{mm/damon/} \newline
\kpath{mm/hmm.c} \newline
\kpath{mm/memory-tiers.c} \\
\bottomrule
\end{tabularx}
\end{frame}

\begin{frame}[t]{第一批要背下来的关键函数}
\small
\begin{columns}[T]
  \begin{column}{0.49\textwidth}
    \begin{block}{分配 / 释放}
    \begin{itemize}
      \item \kcode{alloc_pages()}
      \item \kcode{__alloc_pages()}
      \item \kcode{rmqueue()}
      \item \kcode{free_unref_page()}
    \end{itemize}
    \end{block}

    \begin{block}{地址空间 / 缺页}
    \begin{itemize}
      \item \kcode{handle_mm_fault()}
      \item \kcode{do_anonymous_page()}
      \item \kcode{do_wp_page()}
      \item \kcode{finish_fault()}
    \end{itemize}
    \end{block}

    \begin{block}{页缓存}
    \begin{itemize}
      \item \kcode{filemap_fault()}
      \item \kcode{filemap_get_pages()}
      \item \kcode{page_cache_ra_unbounded()}
      \item \kcode{folio_mark_accessed()}
    \end{itemize}
    \end{block}
  \end{column}
  \begin{column}{0.49\textwidth}
    \begin{block}{回收 / 交换}
    \begin{itemize}
      \item \kcode{shrink_node()}
      \item \kcode{shrink_folio_list()}
      \item \kcode{try_to_unmap()}
      \item \kcode{swap_read_folio()}
      \item \kcode{swap_writeout()}
    \end{itemize}
    \end{block}

    \begin{block}{现代特性}
    \begin{itemize}
      \item \kcode{khugepaged}
      \item \kcode{lru_gen_*}
      \item \kcode{damon_*}
      \item \kcode{hmm_range_fault()}
    \end{itemize}
    \end{block}

    \begin{block}{读函数时的习惯}
    \begin{itemize}
      \item 先找调用方向，再找锁，再看数据结构。
      \item 先认清它处理的是 anon、file、swap、folio 还是 page。
      \item 不要一上来就钻细节分支；先找主路径。
    \end{itemize}
    \end{block}
  \end{column}
\end{columns}
\end{frame}

\begin{frame}[t]{读旧资料时最容易卡住的“术语断层”}
\footnotesize
\begin{tabularx}{\textwidth}{p{0.22\textwidth}p{0.25\textwidth}Y}
\toprule
旧资料里常见说法 & 当前主线更常见说法 & 你应该怎么理解 \\
\midrule
\kcode{bootmem} &
\kcode{memblock} &
很多旧书和老文章先讲 \kcode{bootmem}；今天读当前主线，应把“早期启动内存管理”优先映射到 \kcode{memblock}。 \\
\addlinespace
\kcode{struct page} 作为一切文件缓存对象 &
\kcode{folio} 在 file mm / page cache 路径中越来越主导 &
旧资料讲 page cache 往往全部用 \kcode{page}；现在很多 file mm 路径应优先按 \kcode{folio} 心智模型来读。 \\
\addlinespace
VMA 红黑树 / \kcode{mm->mm_rb} &
\kcode{mm\_struct} 中使用 Maple Tree &
旧资料讲“查找 VMA 主要靠 rb-tree”；当前主线读 \kpath{Documentation/mm/process_addrs.rst} 时应转成 Maple Tree 视角。 \\
\addlinespace
active / inactive LRU 就是 reclaim 的全部故事 &
当前主线可能启用 multigen LRU，旧模型仍有帮助但不是唯一解释 &
旧 reclaim 文章仍值得读，但要知道现代回收路径有新观测接口和新策略。 \\
\addlinespace
高端内存 / \kcode{highmem} 是核心主题 &
在现代 64 位主线上，重点更多转向 folio、reclaim、memcg、heterogeneous memory &
高端内存仍是历史知识点，但不应占用新手前期过多精力。 \\
\bottomrule
\end{tabularx}
\end{frame}

\begin{frame}[t]{版本演进总览：先抓“理解主线”}
\footnotesize
\begin{tabularx}{\textwidth}{p{0.19\textwidth}p{0.22\textwidth}Y}
\toprule
时代 & 代表性变化 & 对新手理解的意义 \\
\midrule
Linux 0.x / 1.x &
从 386 分页世界出发，\kcode{swapper_pg_dir}、基本 buddy、基本 demand paging &
把 \texttt{mm} 先看成“页表 + 页框分配器 + 简单缺页”。 \\
\addlinespace
2.0 / 2.2 / 2.4 &
更大内存、更强 SMP 压力，zone / highmem / page cache 重要性上升 &
开始形成“物理内存布局 + 虚拟地址空间 + 缓存与回收”三件事一起看。 \\
\addlinespace
2.6 大时代 &
rmap、NUMA、memcg、KSM、compaction、migration、hotplug、slab 家族扩展 &
\texttt{mm} 从“分页实现”成长为“系统级内存平台”。 \\
\addlinespace
2.6.38 到 4.x &
THP、workingset/refault、zswap、5-level page tables 等性能与可扩展性主题 &
开始必须把性能、回收策略、TLB、碎片整理一起理解。 \\
\addlinespace
5.x 到 6.x &
HMM、DAMON、folio、Maple Tree、multigen LRU 等现代化重构 &
新 API、新数据结构、新可观测性同时出现，旧资料需要“翻译”后再读。 \\
\addlinespace
当前主线 7.0-rc7 &
folio 语义更普遍，Maple Tree 已是地址空间基础设施，回收与异构内存继续演进 &
读当前代码时必须优先采用现代心智模型，而不是照搬 2.6 时代叙述。 \\
\bottomrule
\end{tabularx}
\end{frame}

\begin{frame}[t]{阶段 1：Linux 0.x / 1.x 的 \texttt{mm} 该怎么看}
\small
\begin{itemize}
  \item 这是“最小可理解模型”的时代：页表、页框、缺页、中断、换入换出。
  \item 当前文档 \kpath{Documentation/mm/page_tables.rst} 仍明确提到：
  早期 Linux 1.0 使用单个顶层页表 \kcode{swapper_pg_dir}，覆盖的是当时非常典型的 4MB 物理内存世界。
  \item 对新手最有价值的地方不是去背每个历史实现细节，而是意识到：
  \begin{itemize}
    \item 起点非常 x86 / 386 导向。
    \item 页表层级比今天浅得多。
    \item 当时 \texttt{mm} 规模远小于今天，主问题也更单纯。
  \end{itemize}
  \item 学历史版本的正确姿势：
  \begin{itemize}
    \item 先学“最小模型”为何成立。
    \item 再看今天为什么必须引入更多层：NUMA、memcg、THP、device memory、Maple Tree、MGLRU。
  \end{itemize}
\end{itemize}
\end{frame}

\begin{frame}[t]{阶段 2：2.0 / 2.2 / 2.4，\texttt{mm} 开始面对“真实规模”}
\small
\begin{itemize}
  \item 更大的 RAM、更复杂的 SMP、更多文件缓存与 I/O 压力，让 \texttt{mm} 不再只是“页表实现”。
  \item 这一阶段你应该重点抓住三条历史线索：
  \begin{enumerate}
    \item \textbf{zone / highmem 视角出现并重要化}：物理内存不再被视作单一池子。
    \item \textbf{page cache 重要性急剧上升}：文件 I/O、mmap、缓存重用开始成为核心主题。
    \item \textbf{回收问题变成系统问题}：什么时候回收、回收谁、代价是什么，不再是边角问题。
  \end{enumerate}
  \item 如果你读到老资料大量讨论 \kcode{ZONE\_DMA}、\kcode{ZONE\_HIGHMEM}、buffer/page cache 分工，不要惊讶；
  这些都是从“更大机器和更复杂 I/O”逼出来的。
  \item 今天学这一段历史的意义是：理解为什么现代 \texttt{mm} 永远在平衡“可分配性、可访问性、可回收性、可迁移性”。
\end{itemize}
\end{frame}

\begin{frame}[t]{阶段 3：2.6 是 \texttt{mm} 真正“长成体系”的时代}
\small
\begin{itemize}
  \item 如果说早期版本是“能工作”，那么 2.6 系列开始更像“要在各种规模机器上都能工作”。
  \item 新手应把这段时代记成六个关键词：
  \begin{itemize}
    \item \textbf{rmap}：反向映射让回收、迁移、拆页、失效同步变得更系统。
    \item \textbf{NUMA}：node-local 分配、策略、拓扑意识成为常规主题。
    \item \textbf{memcg}：内存不再只按整机看，也要按 cgroup 控制和统计。
    \item \textbf{compaction / migration}：高阶分配和大页需要对抗碎片。
    \item \textbf{KSM / hotplug / hugetlb}：\texttt{mm} 开始服务虚拟化、大规模部署、动态硬件。
    \item \textbf{slab 家族演进}：内核对象分配成为独立优化主题。
  \end{itemize}
  \item 读这段历史的最佳态度：不要追求“一次看懂 2.6 全部”，而是认识到现代 \texttt{mm} 的大多数支柱都在这时成型。
\end{itemize}
\end{frame}

\begin{frame}[t]{阶段 4：2.6.38 到 4.x，性能模型开始主导设计}
\small
\begin{itemize}
  \item \textbf{THP} 把“大页是否值得、如何整理碎片、何时折叠/拆分”推到前台。
  \item \textbf{workingset / refault} 让 reclaim 不再只是静态 LRU，而是更关注访问历史和“错误驱逐”的代价。
  \item \textbf{zswap / frontswap 一类技术} 说明 swap 也开始从“要不要用”变成“如何减轻代价”。
  \item \textbf{更深页表层级} 说明虚拟地址空间和平台规模继续增长，旧的简单页表想象已不足够。
  \item 对学习者最关键的变化：
  \begin{itemize}
    \item TLB、页表层级、huge page、内存碎片、直接回收和后台回收开始彼此强耦合。
    \item 你不能只会看“功能正确”，还要会看“路径上的性能代价”。
  \end{itemize}
\end{itemize}
\end{frame}

\begin{frame}[t]{阶段 5：5.x / 6.x 到当前主线，现代化重构全面出现}
\small
\begin{itemize}
  \item \textbf{HMM / \kcode{ZONE\_DEVICE}}：\texttt{mm} 不再只管传统 DRAM，还要把 GPU / device memory 融入统一内存语义。
  \item \textbf{DAMON}：访问模式观测从“只靠传统统计”转向“更细粒度、可编排的监测框架”。
  \item \textbf{folio（5.16 前后开始成为关键词）}：file mm / page cache 路径中的基本对象表达方式发生明显变化。
  \item \textbf{Maple Tree（6.1 前后成为地址空间主视角）}：VMA 管理不应再用老的红黑树心智模型硬套。
  \item \textbf{multigen LRU（5.18 前后进入主线并持续成熟）}：现代 reclaim 解释需要新的代际视角。
  \item 对新手的结论只有一句话：
  旧资料仍然有价值，但你必须在脑中完成“旧术语 $\rightarrow$ 当前主线实现”的翻译。
\end{itemize}
\end{frame}

\begin{frame}[t]{当前这棵 7.0-rc7 树里，主导思想是什么}
\small
\begin{block}{你现在读代码时应该优先带着的心智模型}
\begin{itemize}
  \item \textbf{对象层面}：file mm 优先想到 \kcode{folio}，不是旧式“什么都叫 page”。
  \item \textbf{地址空间层面}：VMA 已是 Maple Tree 语义世界，不是旧书里的 \kcode{mm\_rb} 世界。
  \item \textbf{回收层面}：active/inactive LRU 仍是背景知识，但现代主线还要理解 multigen LRU。
  \item \textbf{控制层面}：memcg、NUMA、sysfs、debugfs、tracepoint、page owner 等可观测性非常重要。
  \item \textbf{硬件层面}：\texttt{mm} 需要同时面对 DRAM、huge pages、hotplug、device memory、tiered memory。
\end{itemize}
\end{block}

\begin{block}{一句话总结}
今天的 \texttt{mm} 已经不是“分页子系统”，而是“统一协调分配、映射、缓存、回收、隔离、监控与异构内存的基础平台”。
\end{block}
\end{frame}

\begin{frame}[t]{学习总策略：不要一上来就淹死在细节里}
\small
\begin{enumerate}
  \item 先建立总地图：node / zone / folio / VMA / page tables / reclaim。
  \item 再抓 5 条主路径：分配、缺页、页缓存、回收、控制。
  \item 每学一段都做实验，不做“纯阅读型学习”。
  \item 旧资料只用来解释“为什么会变成今天这样”，不要拿旧术语直接套当前代码。
  \item 始终把文档、代码、运行时观测三件事绑在一起。
\end{enumerate}

\vspace{0.6em}
\begin{block}{建议周期}
完整从零到“可以独立读主线路径并做小改动”，按兼职学习强度建议 \textbf{16 周}；如果全职投入，可压缩到 8 到 10 周，但不建议再短。
\end{block}
\end{frame}

\begin{frame}[t]{准备阶段：第 0 周先搭好学习环境}
\footnotesize
\begin{tabularx}{\textwidth}{p{0.16\textwidth}YY}
\toprule
任务 & 要做什么 & 产出标准 \\
\midrule
源码导航 &
学会 \kcode{rg}、\kcode{git grep}、ctags/cscope/LSP 的最小用法 &
能在 1 分钟内定位 \kcode{alloc_pages}、\kcode{handle_mm_fault}、\kcode{filemap_fault}、\kcode{shrink_node}。 \\
\addlinespace
内核配置 &
准备一个带调试项的配置：\kcode{CONFIG_DEBUG_VM}、\kcode{CONFIG_PAGE_OWNER}、\kcode{CONFIG_PAGE_TABLE_CHECK}、\kcode{CONFIG_MEMCG}、\kcode{CONFIG_LRU_GEN}、\kcode{CONFIG_TRANSPARENT_HUGEPAGE} &
你能解释每个选项是为了观测什么。 \\
\addlinespace
运行环境 &
最好准备 qemu/kvm 或可重启测试机，避免只在生产机观察 &
有一套能安全触发内存压力、交换、THP 的实验环境。 \\
\addlinespace
基础概念 &
先读 \kpath{Documentation/admin-guide/mm/concepts.rst}、\kpath{Documentation/mm/page_tables.rst}、\kpath{Documentation/mm/physical_memory.rst} &
能够口头解释 PFN、TLB、zone、node、anon、page cache、reclaim。 \\
\bottomrule
\end{tabularx}
\end{frame}

\begin{frame}[t]{第 1--2 周：补硬件与页表基础}
\footnotesize
\begin{tabularx}{\textwidth}{p{0.18\textwidth}YY}
\toprule
维度 & 内容 & 最终标准 \\
\midrule
核心目标 &
把“虚拟地址 $\rightarrow$ 页表 $\rightarrow$ 物理页框 $\rightarrow$ TLB”的链条补齐 &
不再把页表、缺页、权限位、TLB 当成抽象口号。 \\
\addlinespace
必读文档 &
\kpath{Documentation/admin-guide/mm/concepts.rst} \newline
\kpath{Documentation/mm/page_tables.rst} &
知道五级页表是当前软件层次表达；知道哪些层可能 fold。 \\
\addlinespace
关键代码 &
\kpath{include/linux/pgtable.h} \newline
\kpath{include/linux/mm_types.h} \newline
\kpath{arch/x86/mm/fault.c}（若以 x86\_64 为主） &
能区分架构相关 fault 入口和通用 \texttt{mm} 路径。 \\
\addlinespace
实验 &
跟一次用户态缺页，记录“地址、VMA、PTE、folio/page”的变化 &
能画出匿名页首次访问和写时复制的基本时序图。 \\
\bottomrule
\end{tabularx}
\end{frame}

\begin{frame}[t]{第 3--4 周：物理内存、zone、buddy、PCP}
\footnotesize
\begin{tabularx}{\textwidth}{p{0.18\textwidth}YY}
\toprule
维度 & 内容 & 最终标准 \\
\midrule
核心目标 &
搞清 Linux 如何把物理内存表示为 node / zone / free area / order / PCP &
能从 \kcode{gfp} 标志一路解释到“最终从哪个 zone 取页”。 \\
\addlinespace
必读文档 &
\kpath{Documentation/mm/physical_memory.rst} \newline
\kpath{Documentation/mm/page_allocation.rst} &
理解 watermarks、zonelist、movable zone、大页分配失败为何常与碎片有关。 \\
\addlinespace
关键代码 &
\kpath{mm/page_alloc.c} \newline
\kpath{include/linux/mmzone.h} \newline
\kpath{include/linux/gfp_types.h} &
能读主分配路径，不被所有 fastpath / slowpath 细节立刻淹没。 \\
\addlinespace
实验 &
看 \kcode{/proc/zoneinfo}、\kcode{/proc/buddyinfo}、\kcode{/proc/pagetypeinfo}，观察压力前后变化 &
能把“高阶分配失败”与“碎片、watermark、fallback”联系起来。 \\
\bottomrule
\end{tabularx}
\end{frame}

\begin{frame}[t]{第 5--6 周：\texttt{slab/slub}、\texttt{vmalloc}、早期内存}
\footnotesize
\begin{tabularx}{\textwidth}{p{0.18\textwidth}YY}
\toprule
维度 & 内容 & 最终标准 \\
\midrule
核心目标 &
区分“页分配器”和“内核对象分配器”；理解早期启动期与运行期分配方式不同 &
不再把 \kcode{kmalloc}、\kcode{alloc_pages}、\kcode{vmalloc} 混为一谈。 \\
\addlinespace
必读文档 &
\kpath{Documentation/mm/slab.rst} \newline
\kpath{Documentation/admin-guide/mm/slab.rst} \newline
\kpath{Documentation/mm/vmalloc.rst} &
知道对象分配、线性映射之外的虚拟映射、cache 热路径分别在解决什么问题。 \\
\addlinespace
关键代码 &
\kpath{mm/slub.c} \newline
\kpath{mm/slab_common.c} \newline
\kpath{mm/vmalloc.c} \newline
\kpath{mm/memblock.c} \newline
\kpath{mm/mm_init.c} &
知道旧资料里的 \kcode{bootmem} 今天如何映射为 \kcode{memblock} 世界。 \\
\addlinespace
实验 &
读 \kcode{/proc/slabinfo}，观察常见 cache；触发不同尺寸 \kcode{kmalloc}；区分 \kcode{vmalloc} 和线性映射地址 &
能解释为什么“连续虚拟地址”不等于“连续物理地址”。 \\
\bottomrule
\end{tabularx}
\end{frame}

\begin{frame}[t]{第 7--8 周：\texttt{mm\_struct}、VMA、page fault、COW}
\footnotesize
\begin{tabularx}{\textwidth}{p{0.18\textwidth}YY}
\toprule
维度 & 内容 & 最终标准 \\
\midrule
核心目标 &
学会从进程地址空间视角读 \texttt{mm}：\kcode{mm_struct}、VMA、页表和 fault 是怎样连起来的 &
能独立讲清匿名缺页、文件缺页、写时复制的主路径。 \\
\addlinespace
必读文档 &
\kpath{Documentation/mm/process_addrs.rst} \newline
\kpath{Documentation/mm/page_tables.rst} &
知道当前主线里 VMA 查询/组织要按 Maple Tree 心智模型理解。 \\
\addlinespace
关键代码 &
\kpath{mm/mmap.c} \newline
\kpath{mm/vma.c} \newline
\kpath{mm/memory.c} \newline
\kpath{mm/mprotect.c} \newline
\kpath{mm/rmap.c} &
能在源码里顺着 fault 主路径走到 PTE 建立、COW、rmap。 \\
\addlinespace
实验 &
用 \kcode{mmap()}、\kcode{fork()}、写入触发 COW；观察 \kcode{/proc/<pid>/maps}、\kcode{smaps}、\kcode{pagemap} &
能把 VMA、页表项、物理页和进程行为对应起来。 \\
\bottomrule
\end{tabularx}
\end{frame}

\begin{frame}[t]{第 9--10 周：page cache、readahead、writeback、folio}
\footnotesize
\begin{tabularx}{\textwidth}{p{0.18\textwidth}YY}
\toprule
维度 & 内容 & 最终标准 \\
\midrule
核心目标 &
把“文件 I/O 为什么会进入 \texttt{mm}”这件事真正搞懂 &
知道 page cache 为何是 Linux 文件路径和内存路径之间的核心桥梁。 \\
\addlinespace
必读文档 &
\kpath{Documentation/mm/page_cache.rst} \newline
\kpath{Documentation/mm/process_addrs.rst}（file-backed VMA 部分） &
理解当前文档已经明确把 folio 作为 page cache 的基本管理单位。 \\
\addlinespace
关键代码 &
\kpath{mm/filemap.c} \newline
\kpath{mm/readahead.c} \newline
\kpath{mm/truncate.c} \newline
\kpath{mm/shmem.c} &
能区分 buffered I/O、mmap file fault、readahead、dirty/writeback。 \\
\addlinespace
实验 &
对比 \kcode{read()} 与 \kcode{mmap()} 读文件；观察缓存命中前后、预读和回收行为 &
能解释 folio 为什么能让 file mm 路径表达更清晰。 \\
\bottomrule
\end{tabularx}
\end{frame}

\begin{frame}[t]{第 11--12 周：reclaim、swap、OOM、workingset}
\footnotesize
\begin{tabularx}{\textwidth}{p{0.18\textwidth}YY}
\toprule
维度 & 内容 & 最终标准 \\
\midrule
核心目标 &
理解“内存不够时到底发生了什么”，包括后台回收、直接回收、交换和 OOM &
能把 \kcode{kswapd}、direct reclaim、swap in/out、OOM 串成完整故事。 \\
\addlinespace
必读文档 &
\kpath{Documentation/mm/page_reclaim.rst} \newline
\kpath{Documentation/mm/swap.rst} \newline
\kpath{Documentation/mm/oom.rst} \newline
\kpath{Documentation/admin-guide/mm/multigen_lru.rst} &
知道旧 LRU 叙述与当前 multigen LRU 观测接口之间的关系。 \\
\addlinespace
关键代码 &
\kpath{mm/vmscan.c} \newline
\kpath{mm/workingset.c} \newline
\kpath{mm/swap.c} \newline
\kpath{mm/swapfile.c} \newline
\kpath{mm/oom_kill.c} \newline
\kpath{mm/zswap.c} &
能在代码中找到“为什么被回收”“为什么被换出”“为什么最后触发 OOM”。 \\
\addlinespace
实验 &
人为制造内存压力，观察 \kcode{/proc/vmstat}、\kcode{/sys/kernel/mm/lru_gen/}、swap 活动和 OOM 日志 &
能分辨“抖动”“正常回收”“已经进 OOM 边缘”三种状态。 \\
\bottomrule
\end{tabularx}
\end{frame}

\begin{frame}[t]{第 13--14 周：THP、compaction、migration、NUMA、memcg}
\footnotesize
\begin{tabularx}{\textwidth}{p{0.18\textwidth}YY}
\toprule
维度 & 内容 & 最终标准 \\
\midrule
核心目标 &
进入“真正像现代 \texttt{mm}”的主题：大页、碎片整理、跨节点、分组控制 &
能把性能、可分配性、拓扑、隔离这四件事放在一起思考。 \\
\addlinespace
必读文档 &
\kpath{Documentation/mm/transhuge.rst} \newline
\kpath{Documentation/admin-guide/mm/transhuge.rst} \newline
\kpath{Documentation/mm/page_migration.rst} \newline
\kpath{Documentation/mm/numa.rst} \newline
\kpath{Documentation/admin-guide/mm/numa_memory_policy.rst} &
知道 THP、mTHP、compaction、migration、NUMA policy、memcg 各自解决什么问题。 \\
\addlinespace
关键代码 &
\kpath{mm/huge_memory.c} \newline
\kpath{mm/khugepaged.c} \newline
\kpath{mm/compaction.c} \newline
\kpath{mm/migrate.c} \newline
\kpath{mm/mempolicy.c} \newline
\kpath{mm/memcontrol.c} &
能顺着 THP 折叠或 compaction 失败路径读出约束条件。 \\
\addlinespace
实验 &
启用/关闭 THP，不同 \kcode{madvise} 策略；观察 \kcode{khugepaged}、碎片整理、NUMA 绑定和 memcg 压力 &
能解释“为什么某些 workload 适合 THP，某些不适合”。 \\
\bottomrule
\end{tabularx}
\end{frame}

\begin{frame}[t]{第 15--16 周：现代主题收尾，建立“当前主线感”}
\footnotesize
\begin{tabularx}{\textwidth}{p{0.18\textwidth}YY}
\toprule
维度 & 内容 & 最终标准 \\
\midrule
核心目标 &
把 modern mm 的新话语体系真正装进脑子：folio、Maple Tree、MGLRU、DAMON、HMM、memory tiers &
能看懂近几年的主线 \texttt{mm} 讨论，不再只会引用老书。 \\
\addlinespace
必读文档 &
\kpath{Documentation/mm/process_addrs.rst} \newline
\kpath{Documentation/mm/multigen_lru.rst} \newline
\kpath{Documentation/mm/damon/design.rst} \newline
\kpath{Documentation/mm/hmm.rst} &
知道当前 \texttt{mm} 为什么越来越强调可观测性、异构内存、现代数据结构。 \\
\addlinespace
关键代码 &
\kpath{mm/vma.c} \newline
\kpath{mm/damon/} \newline
\kpath{mm/hmm.c} \newline
\kpath{mm/memory-tiers.c} &
能指出这些主题分别作用在“地址空间”“回收策略”“异构设备”“冷热分层”的哪一层。 \\
\addlinespace
最终项目 &
做一个 mini project：写学习笔记、画主路径图、加 tracepoint/bpftrace 观测脚本，或修一个小 warning / 注释 / 文档问题 &
达到“能开始参与 \texttt{mm} 社区讨论”的最低门槛。 \\
\bottomrule
\end{tabularx}
\end{frame}

\begin{frame}[t]{每周固定节奏：不要只读不练}
\small
\begin{enumerate}
  \item 周一：读文档，补概念，写 1 页笔记。
  \item 周二：从 1 个关键函数出发跟主调用链，最多追两层分支。
  \item 周三：做 1 个实验，必须有观测结果，不接受“脑补运行”。
  \item 周四：回到代码解释实验现象，补齐锁、数据结构、状态转换。
  \item 周五：输出 1 张时序图或 1 张数据结构图。
  \item 周末：复盘，把“旧术语”和“当前实现”做一次映射。
\end{enumerate}

\vspace{0.6em}
\begin{block}{最重要的纪律}
每学完一个主题，必须能回答三个问题：
\begin{itemize}
  \item 这个机制处理的核心对象是什么？
  \item 它的快路径和慢路径分别在哪里？
  \item 它和 reclaim / fault / folio / memcg / NUMA 中的哪条线有耦合？
\end{itemize}
\end{block}
\end{frame}

\begin{frame}[t]{必须做的 8 个实验}
\small
\begin{enumerate}
  \item 匿名页首次访问与 COW：观察 \kcode{fork()} 前后写入行为。
  \item zone / buddy 观测：在低压和高压下对比 \kcode{/proc/buddyinfo}。
  \item \kcode{kmalloc} vs \kcode{alloc_pages} vs \kcode{vmalloc}：区分对象、页、虚拟映射。
  \item buffered I/O vs \kcode{mmap()}：观察 page cache 与 file fault。
  \item direct reclaim vs \kcode{kswapd}：制造不同级别压力，观察 stall。
  \item THP / mTHP：切换 \kcode{always/madvise/never}，看性能和内存碎片表现。
  \item memcg 压力：在 cgroup 中跑 workload，看回收与 OOM 是否局部化。
  \item multigen LRU：查看 \kcode{/sys/kernel/mm/lru_gen/} 与 debug 接口，理解冷热代际分布。
\end{enumerate}
\end{frame}

\begin{frame}[t]{观测与调试工具箱}
\footnotesize
\begin{tabularx}{\textwidth}{p{0.18\textwidth}YY}
\toprule
工具 / 接口 & 用来观察什么 & 何时最有用 \\
\midrule
\kcode{/proc/meminfo}、\kcode{/proc/vmstat} &
整机内存概况、换页、回收、脏页、OOM 前兆 &
刚开始做任何内存实验时的第一眼。 \\
\addlinespace
\kcode{/proc/zoneinfo}、\kcode{/proc/buddyinfo}、\kcode{/proc/pagetypeinfo} &
zone 水位、碎片、高阶分配可能性 &
怀疑 compaction / 高阶分配失败时。 \\
\addlinespace
\kcode{/proc/<pid>/maps}、\kcode{smaps}、\kcode{pagemap} &
地址空间、VMA、RSS、匿名/文件页分布 &
研究 fault、COW、mmap、page cache 时。 \\
\addlinespace
\kcode{/sys/kernel/mm/lru_gen/} &
multigen LRU 开关与统计 &
读现代 reclaim 时必看。 \\
\addlinespace
\kcode{tracefs} / \kcode{perf} / \kcode{bpftrace} &
路径级时序与热点 &
确认到底慢在哪一段。 \\
\addlinespace
\kcode{page_owner}、\kcode{page_table_check} &
页来源、非法页表问题 &
查泄漏、页表问题、可疑分配来源。 \\
\addlinespace
\kcode{/proc/slabinfo} &
对象缓存行为 &
看内核对象分配、cache 热度、异常增长。 \\
\bottomrule
\end{tabularx}
\end{frame}

\begin{frame}[t]{读 old mm 资料时，如何做“翻译”而不是照抄}
\small
\begin{block}{正确姿势}
\begin{itemize}
  \item 如果老资料在讲 \kcode{bootmem}：立刻在当前树中对照 \kcode{memblock}。
  \item 如果老资料在讲 page cache 全部是 \kcode{struct page}：去当前 \kpath{Documentation/mm/page_cache.rst} 看 folio 语义。
  \item 如果老资料在讲 VMA 红黑树：回到 \kpath{Documentation/mm/process_addrs.rst}，按 Maple Tree 视角重建理解。
  \item 如果老资料只讲双链表 LRU：再补 \kpath{Documentation/mm/multigen_lru.rst} 和 \kpath{Documentation/admin-guide/mm/multigen_lru.rst}。
  \item 如果老资料完全不提 memcg / HMM / DAMON：把它当“时代背景缺失”，不是你理解错了。
\end{itemize}
\end{block}

\begin{block}{错误姿势}
\begin{itemize}
  \item 把 2.6 时代的术语原样套到 7.0 主线代码上。
  \item 只读经典书，不读当前文档和当前代码。
  \item 把历史设计原因当成当前实现细节。
\end{itemize}
\end{block}
\end{frame}

\begin{frame}[t]{什么时候算“掌握了 \texttt{mm}”}
\small
\begin{enumerate}
  \item 你能解释 node、zone、page、folio、VMA、PTE、anon、page cache、rmap、reclaim、memcg、THP、HMM、DAMON。
  \item 你能独立跟完 4 条主路径：
  \begin{itemize}
    \item 一次页分配
    \item 一次匿名缺页
    \item 一次文件缺页
    \item 一次回收导致的换出或 OOM
  \end{itemize}
  \item 你能把旧资料与当前主线差异说清楚，而不是只会引用旧图。
  \item 你能设计并执行 1 个内存实验，并用代码解释结果。
  \item 你敢改一个小 patch：文档、注释、trace、统计、warning 修复，任意一种。
\end{enumerate}

\vspace{0.6em}
\begin{block}{更现实的判断}
“掌握 \texttt{mm}”不是一次性完成，而是先能读主路径，再选一个方向深钻：
page allocator、reclaim、file mm、THP、memcg、NUMA、device memory 都可以各自深入很多年。
\end{block}
\end{frame}

\begin{frame}[t]{学到最后，建议你选一条“主修线”}
\small
\begin{columns}[T]
  \begin{column}{0.48\textwidth}
    \begin{block}{如果你偏系统基础}
    \begin{itemize}
      \item page allocator
      \item reclaim / vmscan
      \item compaction / migration
      \item memcg / OOM
    \end{itemize}
    \end{block}

    \begin{block}{如果你偏进程地址空间}
    \begin{itemize}
      \item mmap / VMA
      \item fault / COW
      \item GUP / rmap
      \item Maple Tree / locking
    \end{itemize}
    \end{block}
  \end{column}
  \begin{column}{0.48\textwidth}
    \begin{block}{如果你偏文件与缓存}
    \begin{itemize}
      \item filemap / page cache
      \item readahead / writeback
      \item folio 化改造
      \item shmem / tmpfs
    \end{itemize}
    \end{block}

    \begin{block}{如果你偏现代方向}
    \begin{itemize}
      \item THP / mTHP
      \item multigen LRU
      \item DAMON
      \item HMM / device memory / tiers
    \end{itemize}
    \end{block}
  \end{column}
\end{columns}
\end{frame}

\begin{frame}[t]{推荐阅读顺序}
\small
\begin{enumerate}
  \item \kpath{Documentation/admin-guide/mm/concepts.rst}
  \item \kpath{Documentation/mm/page_tables.rst}
  \item \kpath{Documentation/mm/physical_memory.rst}
  \item \kpath{Documentation/mm/process_addrs.rst}
  \item \kpath{Documentation/mm/page_cache.rst}
  \item \kpath{Documentation/mm/transhuge.rst}
  \item \kpath{Documentation/mm/multigen_lru.rst}
  \item \kpath{Documentation/mm/damon/design.rst}
  \item \kpath{Documentation/mm/hmm.rst}
  \item 最后再回头读旧书或老文章，做历史映射
\end{enumerate}

\vspace{0.6em}
\begin{block}{补充建议}
经典旧资料依然有价值，尤其适合建立“为什么当年要这样设计”的感觉；但只要你开始读当前主线代码，就必须优先相信当前文档和当前源码。
\end{block}
\end{frame}

\begin{frame}[t]{一句话收尾}
\large
\begin{block}{真正的学习路径}
先把 Linux \texttt{mm} 从“神秘黑箱”变成一张可导航的地图；
再把历史版本变迁当成理解现代设计的背景；
最后通过实验和代码阅读，把这张地图变成你自己的工程直觉。
\end{block}

\vspace{1em}
\normalsize
如果你严格按这份 16 周路线执行，到结束时，你不一定“精通整个 \texttt{mm}”，但已经足够进入主线语境，并开始做真正有价值的深入学习。
\end{frame}

\end{document}
